\documentclass[10pt]{book}

\usepackage[T1]{fontenc}
\usepackage{kpfonts}
\usepackage{amsmath}
\usepackage{amssymb}
\usepackage{amsthm}
\usepackage{mathtools}
\usepackage{empheq}
\usepackage{pgfplots}
\usepackage{cleveref}
\usepackage{graphicx}
\usepackage{geometry}
\usepackage{titlesec}
\usepackage{fancyhdr}
\usepackage{setspace}
\usepackage{microtype}
\usepackage{parskip}
\usepackage{tcolorbox}
\usepackage{booktabs}
\usepackage{float}

\geometry{a5paper, top=1.5cm, bottom=2cm, inner=1.8cm, outer=1.2cm}

\titleformat{\chapter}[display]{\normalfont\large\bfseries}{\chaptertitlename\ \thechapter}{10pt}{\large}
\titleformat{\section}{\normalfont\normalsize\bfseries}{\thesection}{1em}{}
\titleformat{\subsection}{\normalfont\small\bfseries}{\thesubsection}{1em}{}

\pgfplotsset{compat=1.18}

\title{MATH1014 Algebra}
\author{London Wallace}
\date{Semester 2, 2026}

\renewcommand{\chaptermark}[1]{\markboth{#1}{}}

\pagestyle{fancy}
\fancyhf{}
\fancyfoot[C]{\thepage}
\renewcommand{\headrulewidth}{0pt}
\fancypagestyle{plain}{\fancyhf{}\fancyfoot[C]{\thepage}\renewcommand{\headrulewidth}{0pt}}

\fancypagestyle{plain}{\fancyhf{}\renewcommand{\headrulewidth}{0pt}}

\setstretch{1.15}

\titlespacing{\chapter}{0pt}{-10pt}{10pt}

\newcommand{\newthought}[1]{\vspace{0.5\baselineskip}\noindent{\scshape #1}}

\renewcommand{\bibname}{Bibliography and Suggested Reading}

\begin{document}
\frontmatter
\maketitle
\tableofcontents

\chapter{Preface}
\newthought{These notes} follow, and are meant to cement the course \emph{Algebra}, taught at the
University of Western Australia by Dr. Gordon Royle in Semester 2 of 2026. Where needed, the content of
lectures is supplemented by \emph{Linear Algebra Done Right} by Sheldon Axler.

\section{What is Algebra?}

\newthought{I do not yet know enough to say.}

\section{Table of Symbols}
\begin{center}
  \begin{tabular}{cc}
    \toprule
    Symbol          & Meaning                       \\
    \midrule
    $\forall$       & For every \emph{or} for all   \\
    $\exists$       & There exists                  \\
    $ \in $         & In \emph{or} is an element of \\
    $ \subset $     & Is a proper subset of         \\
    $ \subseteq $   & Is a subest of                \\
    $ \cup $        & Union with                    \\
    $ \cap $        & Intersection of               \\
    $ \backslash$   & Not \emph{or} without         \\
    $ \varnothing $ & Empty set                     \\
    $ > $           & Greater than                  \\
    $ < $           & Less than                     \\
    $ \ge $         & Greater than or equal to      \\
    $ \le $         & Less than or equal to         \\
    $|n|$           & The absolute value of $n$     \\
    $\qedsymbol $   & It has been proven            \\
    \bottomrule
  \end{tabular}
\end{center}
\mainmatter

\chapter{Fields}

\section{Lecture 1}

\subsection{Binary Operations}
\newthought{Binary operations} are \emph{operations}\footnote{Operations are not yet defined.} that take two inputs
and produce one output. Addition $(+)$ and multiplication $(\times)$ are common examples of binary operations.

A uniary operation is one that takes only one input, and produces one output; for example, sine. The function
$f(x,y,z)=x+y+z$ is an example of a \emph{ternary} operation. There are an infinite amount of $n$-ary operations, however
binary operations form the basis of the study of fields and rings.

\subsection{Axioms}
\newthought{Axioms} are statements, assumed to be true, used to define the domain of a particular mathematical theory.
The goal is the find the smallest set of obviously true statements that define the domain which we are interested. If any
axioms are logical consequences of others, they should be omitted, and likewise if any properties of the underlying structure
we wish to define are lost, an axiom should be added to its definition to provide those properties. Axioms form a basis for
further reasoning. We cannot prove anything without first agreeing on the postulates of our theories.

\subsection{The Definition of a Field}
\newthought{A field} is a \emph{set} $S$, equipped with the two binary operations $+$ and $\times$, that behave similarly to how
addition and multiplication would behave on real numbers, that also satisfies the following \emph{axioms}:
\newpage
\begin{enumerate}
  \item Commutativity of $+$. $[a+b=b+a]$
  \item Associativity of $+$. $[a+(b+c)=(a+b)+c]$
  \item Existence of the additive identity 0. $[a+0=a]$
  \item Existence of the additive inverse. $[\forall a \in S,\ \exists b \in S : a+b=0]$
  \item Commutativity of $\times$. $[a \times b = b \times a]$
  \item Associativity of $\times$. $[a \times (b \times c) = (a \times b) \times c]$
  \item Existence of the multiplicative identity 1. $[a \times 1 = a]$
  \item Existence of the multiplicative inverse. $[\forall a \in S,\ \exists b \in S : a \times b = 1]$
  \item Distributivity. $[a \times (b + c) = a \times b + a \times c]$
  \item Non-triviality. $[0 \ne 1]$
\end{enumerate}

\subsection{$\mathbb{R}$, $\mathbb{C}$, and $\mathbb{Q}$ are Fields}
\newthought{Starting with the reals}, it is quite clear that all the axioms are met. For any number $a \in \mathbb{R}$, its
additive inverse would be $-a$, and its multiplicative inverse would be $a^{-1}$. The identities represented by the symbols $0$
and $1$ above are in fact $0$ and $1$, and axioms 1, 2, 5, 6, 9, and 10 are obvious. The same is true for the rationals and the
complex numbers.

\subsection{$\mathbb{Z}$ is not a Field}
\newthought{Not every subset} of a field is a field. The integers are not a field because it is not true that for any
$a \in \mathbb{Z},\ \exists b \in \mathbb{Z}$ with $a \times b = 1$. 2 is certainly an integer, but for $2 \times b = 1$ to
hold, $b=\frac12$, which is certainly not an integer!

\subsection{The Binary Field}
\newthought{The set} $\{0,1\}$ with multiplication and addition as usual, except that $1+1\coloneqq 0$ is a field. It is the
smallest field and the unique field with two elements. Addition defined in this way is known in computer science as XOR.

\subsection{The Properties of Fields}
\newthought{From the axioms} there is a standard set of proofs used to establish the properties of fields usually taken for
granted.
\begin{enumerate}
  \item Uniqueness of 0.

    The additive identity 0 is unique.
    Let $0$ and $0'$ be additive identities, then $0+0'=0$ and $0'+0=0'$. From the commutativity of addition, $0+0'=0'+0$, and
    thus $0=0'$. \qed

  \item Uniqueness of 1.

    Let $1$ and $1'$ be multiplicative identities, then $1 \times 1' =1$ and $1' \times 1=1'$. But $1 \times 1' = 1' \times 1$,
    thus $1=1'$ \qed

  \item Uniqueness of the inverses.

    Let $b$ and $c$ be additive inverses of $a$, then
    \begin{align*}
      b &= b + 0 \\
      b &= b + (a+c) \\
      b &= (a+b)+c \\
      b &= 0 + c \\
      b &= c \tag*{\qed}
    \end{align*}

    Let $b$ and $c$ be multiplicative inverses of $a$, then
    \begin{align*}
      b &= b \times 1 \\
      b &= b \times (a \times c) \\
      b &= (a \times b) \times c \\
      b &= 1 \times c \\
      b &= c \tag*{\qed}
    \end{align*}

  \item Cancellation laws

    If $a+b=a+c$, then $b=c$, and if $a \ne 0$ and $ab=ac$ then $b=c$.

  \item Arithmetic of 0.

    $a \times 0 = 0$
  \item Sign rules

    $(-1)\cdot a = -a$ and $(-a)\cdot b = -(ab) = a \cdot (-b)$

  \item No zero divisors

    If $ab=0$ then $a=0$ or $b=0$. A zero divisor is a non zero element $a$, such that $\exists b \ne 0$ with $ab=0$.
\end{enumerate}

\subsection{Power Sets}
\newthought{A power set} $P(S)$ is the set of all subsets of $S$. The set $S=\{1,2,3\}$ has
$P(S)=\{ \{1\}, \{2\}, \{3\}, \{1,2\}, \{1,3\}, \{2,3\}, \{1,2,3\}, \varnothing \}$. If we define multiplication to be
$A \cap  B$ and addition to be $A \cup B\ \backslash A \cap B$, then all the axioms of a field are met \emph{except} the
existence of a multiplicative inverse. The multiplicative identity is the set $\{1,2,3\}$ since the intersection of any
member of $P(S)$ with $\{1,2,3\}$ is itself, but it is not possible for every member of $P(S)$ to intercept with $\{1,2,3\}$
and produce $\{1,2,3\}$. A structure that satisfies every axiom of a field except the existence of a multiplicative inverse
is called a \emph{commutative unital ring}.

\section{Lecture 2}

\end{document}
